\documentclass[a4paper, 12pt]{article}
\usepackage{amsmath}

\newcommand{\templates}{../../template}
\usepackage[a4paper, margin=2.5cm]{geometry}

\usepackage{enumitem}
\setlist[itemize]{noitemsep}
\setlist[enumerate]{noitemsep}

\let\oldpar\paragraph
\renewcommand{\paragraph}[1]{\oldpar{#1\\}\noindent}
\usepackage{graphicx}
\usepackage{hyperref}
\usepackage{makecell}

\newcommand{\settitolo}[1]{\newcommand{\titolo}{#1\\}}
\newcommand{\setprogetto}[1]{\newcommand{\progetto}{#1\\}}
\newcommand{\setcommittenti}[1]{\newcommand{\committenti}{#1\\}}
\newcommand{\setredattori}[1]{\newcommand{\redattori}{#1\\}}
\newcommand{\setrevisori}[1]{\newcommand{\revisori}{#1\\}}
\newcommand{\setresponsabili}[1]{\newcommand{\responsabili}{#1\\}}
\newcommand{\setversione}[1]{
	\ifdefined\versione\renewcommand{\versione}{#1\\}
	\else\newcommand{\versione}{#1\\}\fi
}
\newcommand{\setdestuso}[1]{\newcommand{\uso}{#1\\}}
\newcommand{\setdescrizione}[1]{\newcommand{\descrizione}{#1\\}}

\newcommand{\makefrontpage}{
	\begin{titlepage}
		\begin{center}

		\includegraphics[width=0.4\textwidth]{../../template/WSWS-logos_transparent_crop}\\

		{\Large Winning Software Solution}\\[6pt]
		\href{mailto://winningsoftwaresolution@gmail.com}{winningsoftwaresolution@gmail.com}\\
		
		\ifdefined\progetto
		\vspace{1cm}
		{\Large\progetto}
		{\large\committenti}
		\else\fi
		
		\vspace{1.5cm}
		{\LARGE\titolo}
		
		\vfill
		
		\begin{tabular}{r | l}
		\multicolumn{2}{c}{\textit{Informazioni}}\\
		\hline
		
		\ifdefined\redattori
			\textit{Redattori} &
			\makecell[l]{\redattori}\\
		\else\fi
		\ifdefined\revisori
			\textit{Revisori} &
			\makecell[l]{\revisori}\\
		\else\fi
		\ifdefined\responsabili
			\textit{Respondabili} &
			\makecell[l]{\responsabili}\\
		\else\fi
		
		\ifdefined\versione
			\textit{Versione} & \versione
		\else\fi
		
		\textit{Uso} & \uso
		
		\end{tabular}
		
		\vspace{2cm}
		
		\ifdefined\descrizione
		Descrizione
		\vspace{6pt}
		\hrule
		\descrizione
		\else\fi
		\end{center}
	\end{titlepage}
}
\usepackage{hyperref}
\usepackage{array}
\usepackage{tabularx}

\def\vers#1-#2-#3-#4-#5\\{#1&#2&#3&#4&#5\\\hline}

\newcommand{\addversione}[5]{
	\ifdefined\versioni
		\let\old\versioni
		\renewcommand{\versioni}{#1&#2&#3&#4&#5\\\hline\old}
	\else
		\newcommand{\versioni}{#1&#2&#3&#4&#5\\\hline}
	\fi
}

\newcommand{\setversioni}[1]{\newcommand{\versioni}{#1}}

\newcommand{\makeversioni}{
	\begin{center}
		\begin{tabularx}{\textwidth}{|c|c|c|c|X|}
		\hline
		\textbf{Versione} & \textbf{Data} & \textbf{Persona} & \textbf{Attivtà} & \textbf{Descrizione} \\
		\hline
		\versioni
		\end{tabularx}
	\end{center}
	\clearpage
}

\settitolo{Piano di Qualifica}
\setredattori{Raffaele Oliviero \\ Elia Scandaletti}
\setdestuso{esterno}
\setdescrizione{
Questo documento serve a definire le metriche e i criteri di accettazione dei prodotti.
}

\addversione{0.0.0}{09/01/2021}{Raffaele Oliviero}{Redazione}{Stesura iniziale.}
\addversione{0.0.1}{16/01/2021}{Elia Scandaletti}{Redazione}{Correzione indice di Gulpease.}
\addversione{0.0.2}{16/01/2021}{Giovanni Cocco}{Redazione}{Migliorata la leggibilità.}

\begin{document}

\makefrontpage

\makeversioni

\section{Premessa}
Il seguente documento definisce metriche e requisiti minimi per l'accettazione di una PR.\\
Il documento è suddiviso per categoria di prodotti.\\\\
Queste metriche dovranno essere calcolabili in maniera automatica.
\section{Documenti}
\subsection{Metriche}
\subsubsection{Indice di Gulpease}
L'indice Gulpease è un indice di leggibilità di un testo tarato sulla lingua italiana.\\\\
Ha il vantaggio di utilizzare la lunghezza delle parole in lettere anziché in sillabe.\\
Questo semplifica il calcolo automatico.\\
I risultati sono compresi tra 0 e 100. Il valore "100" indica la leggibilità più alta.\\\\
In generale risulta che testi coi seguenti indici sono difficili da leggere.
\begin{itemize}
    \item < 80 per chi ha la licenza elementare;
    \item < 60 per chi ha la licenza media;
    \item < 40 per chi ha un diploma superiore.
\end{itemize}

\[ \text{Indice di Gulpease} = 89 + \frac{300*\text{numero di frasi} - 10*\text{numero di lettere}}{\text{numero di parole}} \]

I parametri saranno calcolati:
\begin{itemize}
	\item frasi: i punti presenti nel testo;
	\item lettere: i caratteri alfanumerici;
	\item parole: le parole.
\end{itemize}

\subsection{Requisiti minimi}
\subsubsection{Indice di Gulpease}
\begin{itemize}
	\item \textbf{Documenti interni}: minimo 40;
	\item \textbf{Documenti esterni}: minimo 60.
\end{itemize}

\section{Software}

\end{document}